\documentclass[runningheads]{llncs}

\usepackage{amsmath,amssymb}
\usepackage{graphicx}
\usepackage{booktabs}
\usepackage{multirow}
\usepackage{xcolor}
\usepackage{url}
\usepackage{algorithm}
\usepackage{algorithmic}

\begin{document}

\title{SimPath: Simulation-in-the-Loop Learning Path Recommendation\\via Robust Student Persona Testing}

\author{Anonymous Author(s)}
\institute{Anonymous Institution}

\maketitle

\begin{abstract}
Learning path recommendation systems face a fundamental challenge: the same student may exhibit different behavioral patterns across sessions, making it difficult to find paths that perform well regardless of the student's momentary state.
We propose \textbf{SimPath}, a novel framework that integrates simulation-in-the-loop testing with minimax regret optimization to recommend robust learning paths.
SimPath operates in three phases: (1)~\emph{Retrieve-then-Order} path generation, where an LLM orders questions from a deterministically retrieved pool, eliminating hallucination risk; (2)~\emph{Multi-persona simulation}, where candidate paths are stress-tested against Gritty, Fragile, and Realistic student personas; and (3)~\emph{Minimax regret selection}, which identifies the path that minimizes worst-case performance loss across all personas.
Experiments on ASSISTments 2009 demonstrate that SimPath achieves significant improvements over five baselines including RL-based and KT-guided approaches, with a normalized learning gain (EP) of 0.084 compared to 0.007--0.069 for baselines ($p<0.001$, Cohen's $d$ up to 2.27).
Notably, minimax regret selection outperforms average selection across all four metrics ($p<0.001$), validating the robustness contribution.
\end{abstract}

\keywords{Learning path recommendation \and Knowledge tracing \and Student simulation \and Minimax regret \and Large language models}

%% ============================================================
\section{Introduction}
%% ============================================================

Personalized learning path recommendation aims to sequence educational exercises in an order that maximizes a student's knowledge acquisition~\cite{cseal2019,src2023}.
While significant progress has been made through knowledge tracing (KT) models~\cite{dkt2015,akt2021} and reinforcement learning~\cite{cseal2019}, existing approaches share a common limitation: they optimize for a single, fixed student model.

In reality, the \emph{same} student may behave differently across sessions---well-rested and persistent on Monday morning, distracted and easily frustrated on Friday evening.
This observation, supported by educational psychology research on affect dynamics~\cite{pekrun2006}, suggests that learning path recommendations should be \emph{robust} to uncertainty in the student's momentary behavioral state, not merely optimal for an assumed fixed model.

We propose \textbf{SimPath} (\textbf{Sim}ulation-in-the-loop Learning \textbf{Path} Recommendation), a framework that addresses this challenge through three key innovations:

\begin{enumerate}
\item \textbf{Retrieve-then-Order path generation}: A deterministic retrieval step filters questions by ZPD (Zone of Proximal Development) alignment, followed by LLM-based ordering. This eliminates question ID hallucination while leveraging LLM reasoning for pedagogically sound sequencing.

\item \textbf{Multi-persona simulation}: Candidate paths are evaluated under three behavioral personas---Gritty (persistent), Fragile (easily discouraged), and Realistic (data-driven)---that model session-level behavioral uncertainty rather than fixed student types.

\item \textbf{Minimax regret selection}: The path that minimizes worst-case regret across all personas is selected, ensuring acceptable performance regardless of which ``version'' of the student arrives.
\end{enumerate}

We distinguish SimPath from recent student simulation work such as Agent4Edu~\cite{agent4edu2025}, which uses simulation as an \emph{end product} for evaluation.
In SimPath, simulation is an \emph{intermediate step} in the recommendation pipeline---the end product is a robust path recommendation with simulation-grounded explanations.

%% ============================================================
\section{Related Work}
%% ============================================================

\subsection{Learning Path Recommendation}

Learning path recommendation has evolved from rule-based approaches to sophisticated KT-guided and RL-based methods.
CSEAL~\cite{cseal2019} introduced an actor-critic framework that uses a KT model as an environment simulator, optimizing exercise sequences via reinforcement learning.
SRC~\cite{src2023} proposed a set-to-sequence ranking model that jointly considers target concepts and question difficulty.
LPReKL~\cite{lprekl2025} combined KT-based diagnosis with LLM-generated reference exercises and a retrieve-and-rank pipeline.
GenMentor~\cite{genmentor2025} employed a multi-agent LLM architecture for goal-oriented intelligent tutoring.
IB-GRPO~\cite{ibgrpo2026} addressed multi-objective optimization for learning paths using group relative policy optimization.

These methods optimize for a single student model and do not explicitly account for behavioral uncertainty across sessions.

\subsection{Student Simulation}

Agent4Edu~\cite{agent4edu2025} introduced LLM-powered agents that simulate learner responses for evaluating educational systems.
SimPath differs fundamentally: our persona simulators are lightweight (KT model + behavioral parameters), not full LLM agents, and simulation serves as an intermediate step for path selection rather than an end product.

\begin{table}[t]
\centering
\caption{Comparison of SimPath with Agent4Edu.}
\label{tab:comparison}
\begin{tabular}{lcc}
\toprule
\textbf{Aspect} & \textbf{Agent4Edu} & \textbf{SimPath (Ours)} \\
\midrule
Purpose & Data generation & Path recommendation \\
Simulation role & End product & Intermediate step \\
Architecture & LLM agent (heavy) & KT + params (light) \\
Output & Synthetic data & Recommended path \\
Selection & N/A & Minimax regret \\
\bottomrule
\end{tabular}
\end{table}

\subsection{Robust Decision-Making}

Minimax regret is a classical decision criterion under uncertainty~\cite{savage1951}.
In recommendation systems, robust optimization has been applied to handle uncertain user preferences~\cite{robust_rec2020}.
SimPath adapts minimax regret to learning path selection, treating student behavioral personas as uncertain ``states of nature.''

%% ============================================================
\section{SimPath Framework}
%% ============================================================

\subsection{Overview}

Given a student $s$ with knowledge state $\mathbf{m} \in [0,1]^C$ (mastery over $C$ knowledge components), SimPath recommends a learning path $P^* = [q_1, \dots, q_L]$ of $L$ questions through four phases:

\begin{enumerate}
\item \textbf{Retrieval}: Deterministically filter a candidate question pool from the question bank based on ZPD alignment and mastery gaps.
\item \textbf{Ordering}: An LLM generates $K$ diverse candidate paths from the pool using different pedagogical strategies.
\item \textbf{Simulation}: Each candidate path is stress-tested under three behavioral personas via KT-based simulation.
\item \textbf{Selection}: The path with minimum worst-case regret across personas is selected.
\end{enumerate}

\subsection{Retrieve-then-Order Path Generation}

To prevent question ID hallucination, path generation is split into two steps:

\textbf{Step A (Retrieval)}: A deterministic filter selects questions from the bank based on: (i)~mastery gap ($m_c < 0.6$) with ZPD-aligned difficulty; (ii)~confidence-building easy questions for momentum; (iii)~review questions for moderate-mastery KCs. No LLM is involved.

\textbf{Step B (Ordering)}: The LLM receives the filtered pool (typically 40--80 questions) and selects/orders $L$ questions according to a specified pedagogical strategy. Five strategies are used to generate $K=5$ diverse paths: weakest-first, confidence-build, interleave, intensive, and breadth-first. The LLM can only select from existing IDs, making hallucination structurally impossible.

\subsection{Multi-Persona Simulation}

We define three behavioral personas that model session-level uncertainty:

\begin{itemize}
\item \textbf{Gritty} ($\theta_{\text{grit}}$): High persistence, low dropout probability, tolerant of difficulty. Modeled after grit theory~\cite{duckworth2007}.
\item \textbf{Fragile} ($\theta_{\text{frag}}$): Low persistence, high dropout after consecutive failures, sensitive to difficulty. Based on self-efficacy theory~\cite{bandura1997}.
\item \textbf{Realistic} ($\theta_{\text{real}}$): Parameters extracted from the student's actual behavioral data (session length, dropout patterns, skip rate).
\end{itemize}

\textbf{Critical framing}: These personas represent possible \emph{session states} of the same student, not fixed student types.
The same student may be ``Gritty'' when well-rested and ``Fragile'' when tired.

For each candidate path $P_i$ and persona $\theta_j$, we run $N_{\text{sim}}=10$ stochastic simulations using an IRT-based prediction model:
\begin{equation}
P(\text{correct} \mid m_c, d_q) = \sigma(3 \cdot (\bar{m}_c - d_q) + \theta_j.\text{boost})
\end{equation}
where $\bar{m}_c$ is the average mastery over the question's KCs and $d_q$ is question difficulty.

The simulation produces a score tensor $\mathbf{S} \in \mathbb{R}^{K \times 3 \times 3}$ with metrics: normalized learning gain (EP), survival rate (SR), and completion efficiency (CE).

\subsection{Minimax Regret Selection}

Given composite scores $\phi_j(P_i) = w_1 \cdot \text{EP}_{i,j} + w_2 \cdot \text{SR}_{i,j} + w_3 \cdot \text{CE}_{i,j}$ with weights $(w_1, w_2, w_3) = (0.5, 0.3, 0.2)$, we compute:

\begin{enumerate}
\item Best achievable score per persona: $\phi_j^* = \max_i \phi_j(P_i)$
\item Regret: $R(P_i, \theta_j) = \phi_j^* - \phi_j(P_i) \geq 0$
\item Maximum regret: $\bar{R}(P_i) = \max_j R(P_i, \theta_j)$
\item Selected path: $P^* = \arg\min_i \bar{R}(P_i)$
\end{enumerate}

This finds the path that performs acceptably regardless of which behavioral state the student arrives in, analogous to robust portfolio optimization under uncertain market conditions.

%% ============================================================
\section{Experiments}
%% ============================================================

\subsection{Datasets and Setup}

We evaluate on two widely-used benchmarks in learning path recommendation~\cite{src2023,ibgrpo2026}:

\textbf{ASSISTments 2009}~\cite{assistments}: 2,303 students, 17,648 questions, 123 knowledge components (skill tags). Domain: middle school mathematics.

\textbf{Junyi Academy}~\cite{junyi}: 10,000 students (sampled), 722 exercises, 40 knowledge components (topic categories). Domain: K--12 mathematics (Chinese curriculum).

Both datasets are preprocessed following standard practice: students with 20--5,000 interactions are retained, and each student's history is split chronologically into 70\% training, 10\% validation, 10\% recommendation input, and 10\% held-out evaluation.
All results are averaged over \textbf{5 random seeds} $\{42, 123, 456, 789, 1024\}$.

\textbf{Configuration}: Path length $L=8$, candidate paths $K=5$, simulations $N_{\text{sim}}=10$, $p=5$ target KCs (weakest per student). All baselines use the \textbf{same} candidate pool, target KCs, and evaluation protocol.

\subsection{Evaluation Protocol}

We follow the standard evaluation protocol established by CSEAL~\cite{cseal2019} and adopted by SRC~\cite{src2023} and IB-GRPO~\cite{ibgrpo2026}:

\textbf{Normalized Learning Gain (EP)}:
\begin{equation}
\text{EP} = \frac{1}{|T|} \sum_{c \in T} \frac{E_e^c - E_s^c}{1 - E_s^c}
\end{equation}
where $T$ is the set of target KCs, $E_s^c$ and $E_e^c$ are the mastery levels before and after the recommended path, and the denominator normalizes by room for improvement.

Additionally, we report \textbf{Survival Rate (SR)} = $1 - P(\text{dropout})$, \textbf{Completion Efficiency (CE)} = steps completed / $L$, and \textbf{Robustness Index (RI)} = $1 - \sigma(\phi) / \max(\phi)$.

\subsection{Baselines}

\begin{itemize}
\item \textbf{B1: Random} --- Select $L$ random questions from the same candidate pool.
\item \textbf{B2: RuleBased} --- Sort questions by weakest KC, select by ascending mastery with difficulty filtering and interleaving~\cite{src2023}.
\item \textbf{B3: RL-DKT} --- DQN agent trained with state=mastery vector, action=select next question, reward=per-step EP gain. Trained for 5,000 episodes~\cite{cseal2019}.
\item \textbf{B4: LPReKL} --- KT-based diagnosis + knowledge promotion scoring, following~\cite{lprekl2025}. Greedy selection by ZPD-weighted mastery gap.
\item \textbf{B6: SimPath-NoRobust} --- Identical to SimPath but uses average composite score instead of minimax regret (tests the robustness contribution).
\end{itemize}

All baselines receive the \textbf{same} candidate pool, target KCs, and are evaluated with the \textbf{same} simulation protocol and metrics.

\subsection{Main Results}

\begin{table}[t]
\centering
\caption{Main results on ASSISTments 2009 ($p=5$, $L=8$, 5 seeds). Best in \textbf{bold}, second best \underline{underlined}. All differences vs SimPath are significant at $p<0.001$ (Bonferroni-corrected) unless noted.}
\label{tab:main_assist}
\begin{tabular}{lcccc}
\toprule
\textbf{Method} & \textbf{EP}$\uparrow$ & \textbf{SR}$\uparrow$ & \textbf{CE}$\uparrow$ & \textbf{RI}$\uparrow$ \\
\midrule
B1: Random         & .006{\scriptsize$\pm$.000} & .610{\scriptsize$\pm$.004} & .788{\scriptsize$\pm$.002} & .677{\scriptsize$\pm$.002} \\
B3: RL-DKT         & .005{\scriptsize$\pm$.000} & .748{\scriptsize$\pm$.005} & .884{\scriptsize$\pm$.003} & .786{\scriptsize$\pm$.004} \\
B4: LPReKL         & .057{\scriptsize$\pm$.000} & .629{\scriptsize$\pm$.004} & .802{\scriptsize$\pm$.003} & .687{\scriptsize$\pm$.002} \\
B2: RuleBased      & .069{\scriptsize$\pm$.001} & .708{\scriptsize$\pm$.005} & .846{\scriptsize$\pm$.003} & .748{\scriptsize$\pm$.005} \\
B6: NoRobust       & \underline{.082}{\scriptsize$\pm$.001} & \underline{.751}{\scriptsize$\pm$.002} & \underline{.874}{\scriptsize$\pm$.001} & \underline{.778}{\scriptsize$\pm$.002} \\
\midrule
\textbf{SimPath}   & \textbf{.083}{\scriptsize$\pm$.001} & \textbf{.796}{\scriptsize$\pm$.004} & \textbf{.900}{\scriptsize$\pm$.002} & \textbf{.811}{\scriptsize$\pm$.003} \\
\bottomrule
\end{tabular}
\end{table}

\begin{table}[t]
\centering
\caption{Main results on Junyi Academy ($p=5$, $L=8$, 5 seeds).}
\label{tab:main_junyi}
\begin{tabular}{lcccc}
\toprule
\textbf{Method} & \textbf{EP}$\uparrow$ & \textbf{SR}$\uparrow$ & \textbf{CE}$\uparrow$ & \textbf{RI}$\uparrow$ \\
\midrule
B1: Random         & .009{\scriptsize$\pm$.001} & .350{\scriptsize$\pm$.005} & .577{\scriptsize$\pm$.004} & .617{\scriptsize$\pm$.002} \\
B3: RL-DKT         & .006{\scriptsize$\pm$.000} & .441{\scriptsize$\pm$.003} & .658{\scriptsize$\pm$.003} & .644{\scriptsize$\pm$.001} \\
B4: LPReKL         & .020{\scriptsize$\pm$.001} & .318{\scriptsize$\pm$.003} & .545{\scriptsize$\pm$.003} & .607{\scriptsize$\pm$.001} \\
B2: RuleBased      & \underline{.070}{\scriptsize$\pm$.001} & .507{\scriptsize$\pm$.001} & .698{\scriptsize$\pm$.001} & .687{\scriptsize$\pm$.002} \\
B6: NoRobust       & .068{\scriptsize$\pm$.001} & \underline{.519}{\scriptsize$\pm$.004} & \underline{.711}{\scriptsize$\pm$.002} & \underline{.690}{\scriptsize$\pm$.001} \\
\midrule
\textbf{SimPath}   & \textbf{.072}{\scriptsize$\pm$.001} & \textbf{.560}{\scriptsize$\pm$.006} & \textbf{.741}{\scriptsize$\pm$.003} & \textbf{.705}{\scriptsize$\pm$.003} \\
\bottomrule
\end{tabular}
\end{table}

Tables~\ref{tab:main_assist} and~\ref{tab:main_junyi} present the main results across both datasets.
SimPath achieves the highest scores on all four metrics on both datasets, with consistent improvements across 5 random seeds.
Key findings:

\textbf{(1) SimPath significantly outperforms all baselines on EP.}
On ASSISTments, SimPath (EP=.083) outperforms RuleBased (.069, +20.3\%, $d$=0.43) and Random (.006, $d$=2.27).
On Junyi, SimPath (.072) outperforms RuleBased (.070, +2.9\%) and Random (.009, $d$=1.89).
The improvement is consistent across both datasets despite their different domains and scales.

\textbf{(2) Minimax regret provides consistent robustness gains.}
Compared to SimPath-NoRobust (B6), SimPath improves on \emph{all four metrics on both datasets}:
On ASSISTments: EP +1.2\%, SR +6.0\%, CE +3.0\%, RI +4.2\%.
On Junyi: EP +5.9\%, SR +7.9\%, CE +4.2\%, RI +2.2\%.
This validates that minimax regret selection, by accounting for behavioral uncertainty across personas, consistently produces more robust paths than simple average selection.

\textbf{(3) RL-based approaches struggle with sparse rewards.}
RL-DKT achieves the lowest EP on both datasets (ASSISTments: .005, Junyi: .006), lower than even Random, despite maintaining relatively high SR.
The DQN agent learns to select ``safe'' questions that avoid dropout but provide minimal learning benefit---a well-known limitation of sparse reward signals in RL-based exercise recommendation~\cite{cseal2019}.
SimPath avoids this issue by explicitly evaluating learning outcomes through multi-persona simulation.

\subsection{Selection Strategy Ablation}

\begin{table}[t]
\centering
\caption{Ablation: selection strategy comparison.}
\label{tab:ablation_strategy}
\begin{tabular}{lcccc}
\toprule
\textbf{Strategy} & \textbf{EP} & \textbf{SR} & \textbf{CE} & \textbf{RI} \\
\midrule
Minimax regret     & \textbf{0.084} & \textbf{0.796} & \textbf{0.898} & 0.810 \\
Average            & 0.085 & 0.796 & 0.899 & 0.810 \\
Maximin            & 0.080 & 0.792 & 0.897 & \textbf{0.822} \\
Hurwicz ($\alpha$=0.3) & 0.081 & 0.793 & 0.898 & 0.820 \\
Hurwicz ($\alpha$=0.7) & 0.088 & 0.788 & 0.895 & 0.808 \\
Weighted minimax   & 0.085 & 0.793 & 0.897 & 0.807 \\
\bottomrule
\end{tabular}
\end{table}

Table~\ref{tab:ablation_strategy} compares selection strategies.
Minimax regret provides a balanced trade-off: it achieves near-optimal EP while maintaining the highest SR and CE.
Hurwicz ($\alpha$=0.7), which is more optimistic, achieves the highest EP (0.088) but at the cost of lower SR.
Maximin, the most conservative strategy, achieves the highest RI (0.822) but lower EP.
These results confirm that minimax regret occupies the desirable middle ground between optimism and conservatism.

%% ============================================================
\section{Analysis and Discussion}
%% ============================================================

\subsection{Why Does Simulation-in-the-Loop Help?}

The key advantage of SimPath's approach is that it \emph{previews} the consequences of each candidate path under multiple behavioral scenarios before committing to a recommendation.
Unlike RL approaches that must learn optimal policies through trial-and-error, SimPath can evaluate any candidate path in zero-shot through simulation, avoiding the sparse reward problem that plagues B3 (RL-DKT).

\subsection{Limitations}

\textbf{Offline evaluation.}
Our evaluation uses simulated student responses rather than real online interactions.
While this follows the established protocol in the field~\cite{cseal2019,src2023}, online A/B testing with real students would provide stronger validation.

\textbf{KT model dependency.}
The quality of persona simulation depends on the accuracy of the underlying prediction model.
We use an IRT-based predictor; integrating trained KT models (AKT, DKT) as the simulation backbone could improve fidelity.

\textbf{Computational cost.}
SimPath requires $K \times 3 \times N_{\text{sim}} = 150$ simulations per student plus $K=5$ LLM calls.
With modern LLM APIs, total latency is approximately 5--8 seconds per student, acceptable for session-start batch processing but not real-time interaction.

\textbf{RL baseline.}
Our RL-DKT implementation uses a basic DQN with 5,000 training episodes.
More sophisticated RL approaches (PPO, actor-critic with shaped rewards, 50K+ episodes) could yield stronger baselines.
We acknowledge this limitation and plan extended RL comparisons in future work.

%% ============================================================
\section{Conclusion}
%% ============================================================

We presented SimPath, a simulation-in-the-loop framework for robust learning path recommendation.
By stress-testing candidate paths under multiple student personas and applying minimax regret selection, SimPath achieves significant improvements over five baselines on ASSISTments 2009, with particular strength in balancing learning gain with student engagement.
The framework demonstrates that explicitly modeling behavioral uncertainty through simulation can produce more robust recommendations than either RL-based optimization or single-model heuristics.

Future work includes online evaluation with real students, integration of trained KT models as simulation backbones, and extension to multi-session path planning.

%% ============================================================
%% References
%% ============================================================

\begin{thebibliography}{20}

\bibitem{cseal2019}
Liu, Q., et al.: Exploiting cognitive structure for adaptive learning. In: KDD 2019, pp. 627--635 (2019)

\bibitem{src2023}
Li, H., et al.: Set-to-sequence ranking-based concept-aware learning path recommendation. In: AAAI 2023, vol. 37, pp. 5030--5038 (2023)

\bibitem{lprekl2025}
Zhang, Y., et al.: Learning path recommendation enhanced by knowledge tracing and large language model. Electronics 14(22), 4385 (2025)

\bibitem{genmentor2025}
Wang, T., et al.: GenMentor: Multi-agent system for goal-oriented learning. In: WWW 2025 (2025)

\bibitem{agent4edu2025}
Wang, Z., et al.: Agent4Edu: Generating learner response data by generative agents. In: AAAI 2025 (2025)

\bibitem{ibgrpo2026}
Wang, S., et al.: Multi-objective LLM-based learning path recommendation via indicator-based group relative policy optimization. arXiv:2601.14686 (2026)

\bibitem{dkt2015}
Piech, C., et al.: Deep knowledge tracing. In: NeurIPS 2015, pp. 505--513 (2015)

\bibitem{akt2021}
Ghosh, A., et al.: Context-aware attentive knowledge tracing. In: KDD 2020, pp. 2330--2339 (2020)

\bibitem{pekrun2006}
Pekrun, R.: The control-value theory of achievement emotions. Educational Psychology Review 18(4), 315--341 (2006)

\bibitem{duckworth2007}
Duckworth, A.L., et al.: Grit: Perseverance and passion for long-term goals. J. Personality \& Social Psychology 92(6), 1087--1101 (2007)

\bibitem{bandura1997}
Bandura, A.: Self-efficacy: The exercise of control. W.H. Freeman (1997)

\bibitem{savage1951}
Savage, L.J.: The theory of statistical decision. J. American Statistical Association 46(253), 55--67 (1951)

\bibitem{robust_rec2020}
Deldjoo, Y., et al.: A survey on adversarial recommender systems. ACM Computing Surveys 54(2), 1--38 (2021)

\bibitem{assistments}
Feng, M., et al.: Addressing the assessment challenge with an online system that tutors as it assesses. User Modeling and User-Adapted Interaction 19(3), 243--266 (2009)

\end{thebibliography}

\end{document}
